% Figure: the case-study Howe roof truss. Triangular (gable) profile,
% span 16 m in eight bottom-chord panels of 2 m, rise 4 m at the ridge.
% Top chord rises from each support to the central ridge; verticals hang
% from the sloping top chord to the bottom chord; Howe diagonals slope
% up toward the supports. Purlin loads act down at the top-chord joints.
\begin{figure}[htbp]
\centering
\begin{tikzpicture}[>=latex, x=0.78cm, y=0.78cm]
  % bottom chord joints L0..L8 over a 16 m span (1 unit = 1 m)
  \foreach \i in {0,...,8} { \coordinate (L\i) at (\i*2, 0); }
  % ridge height 4 m at midspan; top chord apexes T1..T7 follow the gable
  % gable line: from L0 (0,0) up to ridge (8,4) down to L8 (16,0)
  \coordinate (T1) at (2, 1);
  \coordinate (T2) at (4, 2);
  \coordinate (T3) at (6, 3);
  \coordinate (T4) at (8, 4);   % ridge
  \coordinate (T5) at (10, 3);
  \coordinate (T6) at (12, 2);
  \coordinate (T7) at (14, 1);
  % bottom chord
  \draw[engblue, very thick] (L0) -- (L8);
  % top chord (both rafters)
  \draw[engblue, very thick] (L0) -- (T1) -- (T2) -- (T3) -- (T4)
                                  -- (T5) -- (T6) -- (T7) -- (L8);
  % verticals from each top-chord joint to the bottom chord below it
  \draw[engblue, thick] (T1) -- (L1);
  \draw[engblue, thick] (T2) -- (L2);
  \draw[engblue, thick] (T3) -- (L3);
  \draw[engblue, thick] (T4) -- (L4);
  \draw[engblue, thick] (T5) -- (L5);
  \draw[engblue, thick] (T6) -- (L6);
  \draw[engblue, thick] (T7) -- (L7);
  % Howe diagonals slope up toward the supports (compression diagonals)
  \draw[engblue, thick] (L1) -- (T2);
  \draw[engblue, thick] (L2) -- (T3);
  \draw[engblue, thick] (L3) -- (T4);
  \draw[engblue, thick] (T4) -- (L5);
  \draw[engblue, thick] (T5) -- (L6);
  \draw[engblue, thick] (T6) -- (L7);
  % joints
  \foreach \i in {0,...,8} { \fill[engblack] (L\i) circle (1.5pt); }
  \foreach \p in {T1,T2,T3,T4,T5,T6,T7} { \fill[engblack] (\p) circle (1.5pt); }
  % purlin loads at the seven top-chord joints
  \foreach \p in {T1,T2,T3,T4,T5,T6,T7} {
    \draw[->, very thick, engred] (\p) -- ++(0,-1.0);
  }
  \node[engred, font=\scriptsize, anchor=west] at ($(T4)+(0.1,-0.6)$) {$P$};
  % supports: pin at L0, roller at L8
  \draw[thick] ($(L0)+(-0.3,-0.5)$) -- (L0) -- ($(L0)+(0.3,-0.5)$) -- cycle;
  \draw[thick] ($(L8)+(0,-0.2)$) circle (0.13);
  \draw[thick] ($(L8)+(-0.3,-0.33)$) -- ($(L8)+(0.3,-0.33)$);
  % reactions
  \draw[->, thick, englime!70!black] ($(L0)+(0,-0.6)$) -- ++(0,-0.75);
  \node[font=\scriptsize, anchor=east] at ($(L0)+(-0.05,-1.05)$) {$R_A$};
  \draw[->, thick, englime!70!black] ($(L8)+(0,-0.5)$) -- ++(0,-0.75);
  \node[font=\scriptsize, anchor=west] at ($(L8)+(0.05,-0.95)$) {$R_B$};
  % a few joint labels
  \node[font=\scriptsize, anchor=north east] at (L0) {$L_0$};
  \node[font=\scriptsize, anchor=north] at ($(L1)+(0,-0.02)$) {$L_1$};
  \node[font=\scriptsize, anchor=north] at ($(L4)+(0,-0.02)$) {$L_4$};
  \node[font=\scriptsize, anchor=south] at ($(T2)+(-0.1,0.05)$) {$T_2$};
  \node[font=\scriptsize, anchor=south] at ($(T4)+(0,0.08)$) {$T_4$};
  % span dimension
  \draw[<->, enggray, thin] (0,-1.7) -- (16,-1.7);
  \node[enggray, font=\scriptsize, anchor=north] at (8,-1.7) {$L=16\,\si{\meter}$ (eight $2\,\si{\meter}$ bays)};
  % rise dimension
  \draw[<->, enggray, thin] (16.6,0) -- (16.6,4);
  \node[enggray, font=\scriptsize, anchor=west] at (16.6,2) {rise $4\,\si{\meter}$};
\end{tikzpicture}
\caption[Case-study Howe roof truss]{The case-study Howe roof truss: a
gable profile of span $L=16\,\si{\meter}$ in eight bottom-chord bays of
$2\,\si{\meter}$, rising to a ridge $4\,\si{\meter}$ above the supports.
Purlins deliver equal downward loads $P$ to the seven interior top-chord
joints. The pin at $L_0$ supplies two reaction components and the roller
at $L_8$ supplies one, three reactions against three planar equilibrium
equations, so the truss is externally determinate. The Howe diagonals
slope upward toward the supports, which under gravity load puts them in
compression and the verticals in tension, the reverse of the Pratt
arrangement and the distinction a designer trades on when choosing which
members should be the slender tension elements.}
\label{fig:vol03ch02:howe-roof-truss}
\end{figure}
